% Los 4 modelos del tema 3.A.32 frente a frente, en una cara (comentario de
% Víctor del 19/11/2026): a) IS-LM con OA-DA; b) IS-LM-BP con OA-DA
% (Mundell-Fleming); c) el modelo neokeynesiano canónico; d) la nueva
% macroeconomía abierta, con la pequeña economía abierta (Galí-Monacelli y
% Furlanetto) y los 2 países grandes (Corsetti-Pesenti, con la notación del
% 3.B.19). Filas: supuestos, ecuaciones, qué determina los precios y el tipo
% de cambio, política monetaria, política fiscal y límites.
% Unos 16,6 x 21 cm: cabe en una cara dentro de su caja.
% Autor: Víctor Gutiérrez Marcos
\documentclass[tikz,border=4pt]{standalone}
% ==========================================================================
% tikz-tcee.tex — Estilos comunes de los gráficos TikZ del temario
% Autor: Víctor Gutiérrez Marcos
%
% Cada gráfico es un fichero standalone en fuentes/<tema>/graficos/:
%
%   \documentclass[tikz,border=4pt]{standalone}
%   % ==========================================================================
% tikz-tcee.tex — Estilos comunes de los gráficos TikZ del temario
% Autor: Víctor Gutiérrez Marcos
%
% Cada gráfico es un fichero standalone en fuentes/<tema>/graficos/:
%
%   \documentclass[tikz,border=4pt]{standalone}
%   % ==========================================================================
% tikz-tcee.tex — Estilos comunes de los gráficos TikZ del temario
% Autor: Víctor Gutiérrez Marcos
%
% Cada gráfico es un fichero standalone en fuentes/<tema>/graficos/:
%
%   \documentclass[tikz,border=4pt]{standalone}
%   \input{../../../latex/tikz-tcee}
%   \begin{document}
%   \begin{tikzpicture}[tcee]
%     \ejes{Q}{P}{6}{5}
%     \draw[curva1] (0.5,4.5) -- (5,0.8) node[etiqueta,right] {$D$};
%     \capa{2}{\draw[curva2] ...;}   % en el vídeo aparece en el paso 2
%   \end{tikzpicture}
%   \end{document}
%
% Paleta: la de la web (morado, dorado, salvia) más tres colores de apoyo
% distinguibles en blanco y negro por el trazo.
% ==========================================================================
\usepackage[T1]{fontenc}
\usepackage{newpxtext,newpxmath}
\usepackage{xcolor}
\usepackage{tikz,pgfplots}
\pgfplotsset{compat=1.18}
\usetikzlibrary{arrows.meta,calc,positioning,patterns,decorations.pathreplacing,intersections,shapes.misc,backgrounds}

\definecolor{tceemorado}{HTML}{5F2987}
\definecolor{tceedorado}{HTML}{B8860B}
\definecolor{tceeverde}{HTML}{3C7D3A}
\definecolor{tceeazul}{HTML}{1F5F99}
\definecolor{tceerojo}{HTML}{B22222}
\definecolor{tceenaranja}{HTML}{E07B28}
\definecolor{tceegris}{HTML}{767171}
\definecolor{tceesalvia}{HTML}{E2EFD9}

% Capas para el vídeo: \capa{n}{…} dibuja su contenido solo a partir del paso n.
% En el PDF, la web y Word se ve todo; generar-video.py compila el gráfico una
% vez por paso con \def\tceepaso{k} para que cada elemento aparezca al narrarlo.
% Vectores con flecha encima, como en el texto (tcee.sty): \vect{x}, \vect{p}
\providecommand{\vect}[1]{\vec{#1}}
\providecommand{\tceepaso}{1000}
\newcommand{\capa}[2]{\ifnum#1>\tceepaso\relax\else#2\fi}
% Gráficos por pasos del mismo tamaño (propuesta del 3.A.29, 11/11/2026): con
% \def\tceefantasma{1} antes del \input (p. ej., x-paso1.tex:
% «\def\tceepaso{3}\def\tceefantasma{1}\input{x.tex}»), las capas posteriores
% al paso se dibujan invisibles y ocupan su sitio. Sin \tceefantasma (el vídeo),
% no cambia nada.
\ifdefined\tceefantasma
  \renewcommand{\capa}[2]{\ifnum#1>\tceepaso\relax\begin{scope}[opacity=0]#2\end{scope}\else#2\fi}
\fi

\tikzset{
  tcee/.style={line cap=round,line join=round,font=\small,>={Stealth[length=2.2mm]}},
  eje/.style={->,thick,black},
  curva1/.style={very thick,tceemorado},
  curva2/.style={very thick,tceeazul},
  curva3/.style={very thick,tceeverde},
  curva4/.style={very thick,tceedorado},
  curvanueva/.style={very thick,tceerojo},
  desplazada/.style={very thick,dashed},
  guia/.style={thin,dashed,tceegris},
  punto/.style={circle,fill=black,inner sep=1.3pt},
  etiqueta/.style={font=\small,inner sep=2pt},
  flecha/.style={->,thick,tceerojo},
  area/.style={fill=tceemorado,fill opacity=0.18,draw=none},
  area2/.style={fill=tceedorado,fill opacity=0.22,draw=none},
}

% \ejes{etiqueta x}{etiqueta y}{largo x}{alto y}
\newcommand{\ejes}[4]{%
  \draw[eje] (0,0) -- (#3,0) node[below left=2pt and 0pt] {#1};
  \draw[eje] (0,0) -- (0,#4) node[left] {#2};
  \node[below left] at (0,0) {$0$};}
% \proyeccion{(x,y)}{etq x}{etq y}: líneas guía a los ejes con etiquetas
\newcommand{\proyeccion}[3]{%
  \draw[guia] let \p1=#1 in (\x1,0) node[below] {#2} |- (0,\y1) node[left] {#3};}

% --------------------------------------------------------------------------
% Paleta «dictamen» (dic…): gráficos de datos de la Guía del dictamen (1D01) y
% de cualquier gráfico de coyuntura que quiera los mismos colores. Un concepto,
% un color en todos los gráficos (CRITERIOS §7, «tres reglas», 24/10/2026);
% tabla completa en fuentes/1D01/_trabajo/paleta-1D01.md. No cambia nada de
% lo anterior: son colores y estilos nuevos.
% Actividad: el PIB (y sus variantes: PIB por persona, PIB mundial), morado.
\definecolor{dicpib}{HTML}{5F2987}
% Demanda interna, gama cálida: consumo (naranjas: privado oscuro, público
% claro), inversión (verde) y existencias (gris con trama, aparte).
\definecolor{dicconspriv}{HTML}{C76A1F}
\definecolor{dicconspub}{HTML}{F5C690}
\definecolor{dicfbcf}{HTML}{2A6B3C}
\definecolor{dicexistencias}{HTML}{B4B4B4}
% Demanda externa y comercio, gama fría (azul claro, que se separa del morado)
\definecolor{dicexterna}{HTML}{6AA6DA}
% Zonas: zona del euro (azul marino), España (rojo), resto de países (grises)
\definecolor{dicze}{HTML}{17406D}
\definecolor{dicespana}{HTML}{C8382E}
\definecolor{dicotro}{HTML}{8A8A8A}
\definecolor{dicotroclaro}{HTML}{B5B5B5}
\definecolor{dicotromasclaro}{HTML}{DADADA}
% Precios: inflación general (negro, continuo) y subyacente (negro,
% discontinuo); componentes: energía (bermellón), alimentos (amarillo) y los
% de la subyacente en azules (bienes industriales claro, servicios oscuro).
\definecolor{dicinflacion}{HTML}{1A1A1A}
\definecolor{dicenergia}{HTML}{D55E00}
\definecolor{dicalimentos}{HTML}{F2D45C}
\definecolor{dicbine}{HTML}{7FB2E0}
\definecolor{dicservicios}{HTML}{1F4E89}
% Trabajo: empleo (verde azulado), horas (el mismo, claro), tasa de actividad
% (verde azulado oscuro), paro (teja); productividad (dorado de la web).
\definecolor{dicempleo}{HTML}{1B7F6B}
\definecolor{dichoras}{HTML}{86CBB8}
\definecolor{dicactividad}{HTML}{0F5247}
\definecolor{dicparo}{HTML}{A23B2A}
\definecolor{dicproductividad}{HTML}{B8860B}
% Hogares: renta disponible (pizarra) y tasa de ahorro (negro); el consumo es
% dicconspriv.
\definecolor{dicrenta}{HTML}{4A5A6A}
\definecolor{dicahorro}{HTML}{1A1A1A}
% Política monetaria: tipos del BCE en azules (depósito, el de la zona del
% euro; operaciones principales, medio; facilidad marginal, claro); balance:
% préstamos a los bancos (dorado claro), resto del activo (gris claro).
\definecolor{dicdeposito}{HTML}{17406D}
\definecolor{dicopf}{HTML}{4F86C6}
\definecolor{dicfmc}{HTML}{9EC5EA}
\definecolor{dicprestamos}{HTML}{D9A520}
\definecolor{dicresto}{HTML}{CFCFCF}
% Ciclo: series adelantadas, coincidentes (el PIB) y retrasadas
\definecolor{dicadelantada}{HTML}{D55E00}
\definecolor{dicretrasada}{HTML}{6AA6DA}
% Referencias: rejilla, previsión, sombreados
\definecolor{dicreferencia}{HTML}{767171}

\tikzset{
  % barra de un componente: \addplot[dicbarra=dicfbcf]
  dicbarra/.style={fill=#1, draw=#1!65!black, line width=0.3pt},
  % existencias: gris con trama, para que se distinga también en gris
  dicbarratrama/.style={fill=#1, draw=#1!65!black, line width=0.3pt,
    postaction={pattern=north east lines, pattern color=white}},
  % línea de una serie; con halo blanco cuando pasa sobre barras
  diclinea/.style={very thick, #1},
  diclineahalo/.style={very thick, #1, preaction={draw=white, line width=3.4pt, solid}},
  % puntos de una serie (capa aparte, después de la línea)
  dicpuntos/.style={only marks, mark=*, mark size=1.8pt,
    mark options={fill=white, draw=#1, line width=0.9pt, solid}},
  % línea del cero o de un umbral, después de las barras
  diccero/.style={thin, black},
  dicumbral/.style={thin, dashed, dicreferencia},
}
% Capas de los gráficos de datos (CRITERIOS §7, «tres reglas»), de abajo
% arriba: cuadrícula → sombreados (dicfondo) → barras y áreas (dicbarras) → ejes, cero y umbrales
% (axis lines) → líneas (main) → puntos (dicpuntos) → rótulos (axis
% descriptions). No basta el orden del código: pgfplots dibuja cuadrícula,
% ejes y marcas en sus propias capas (y area style pone los ejes encima), y
% dentro de un axis los \draw, \fill y \node no admiten capa (van con los
% datos de main); por eso todo lo que se dibuja dentro del axis es un \addplot.
% Uso: \begin{axis}[…, dicapas] (al final de las opciones);
%   barras y áreas: \addplot[on layer=dicbarras, …];
%   sombreados:     \dicsombra[estilo]{x1}{y1}{x2}{y2} (rectángulo, en dicfondo,
%                   bajo las barras y sobre la cuadrícula);
%   cero y umbrales: \dicraya[estilo]{(x1,y1) (x2,y2)} (en axis lines);
%   líneas: \addplot normal (main); puntos: \addplot[on layer=dicpuntos, …]
%   (las marcas de cualquier serie van solas a dicpuntos);
%   rótulos: \dicrotulo[opciones del nodo]{(x,y)}{texto} (en axis descriptions),
%   siempre sin fondo (nada de fill=white: si chocan con una línea, se desplazan;
%   CRITERIOS §7, regla 4).
% Los \capa{n} del vídeo no cambian: siguen envolviendo instrucciones completas.
\pgfdeclarelayer{dicfondo}
\pgfdeclarelayer{dicbarras}
\pgfdeclarelayer{dicpuntos}
\pgfplotsset{
  % (el segundo argumento reparte cuadrícula, ejes y rótulos como el estándar)
  layers/dic/.define layer set={axis background, axis grid, dicfondo, dicbarras,
    axis ticks, axis lines, axis tick labels, pre main, main, dicpuntos,
    axis descriptions, axis foreground}{/pgfplots/layers/standard},
  % leyenda sin caja blanca: toma el fondo del gráfico, también si el gráfico va
  % dentro de un cuadro gris (CRITERIOS §7, regla 4, 26/10/2026)
  dicapas/.style={axis on top=false, set layers=dic, mark layer=dicpuntos,
    legend style={fill=none}},
  dicsuelto/.style={forget plot, stack plots=false, sharp plot, mark=none},
}
\newcommand{\dicraya}[2][]{%
  \addplot[on layer=axis lines, dicsuelto, diccero, #1] coordinates {#2};}
\newcommand{\dicsombra}[5][]{%
  \addplot[on layer=dicfondo, dicsuelto, draw=none, #1]
    coordinates {(#2,#3) (#4,#3) (#4,#5) (#2,#5)};}
\newcommand{\dicrotulo}[3][]{%
  \addplot[on layer=axis descriptions, dicsuelto, draw=none] coordinates {#2}
    node[pos=0, #1] {#3};}

%   \begin{document}
%   \begin{tikzpicture}[tcee]
%     \ejes{Q}{P}{6}{5}
%     \draw[curva1] (0.5,4.5) -- (5,0.8) node[etiqueta,right] {$D$};
%     \capa{2}{\draw[curva2] ...;}   % en el vídeo aparece en el paso 2
%   \end{tikzpicture}
%   \end{document}
%
% Paleta: la de la web (morado, dorado, salvia) más tres colores de apoyo
% distinguibles en blanco y negro por el trazo.
% ==========================================================================
\usepackage[T1]{fontenc}
\usepackage{newpxtext,newpxmath}
\usepackage{xcolor}
\usepackage{tikz,pgfplots}
\pgfplotsset{compat=1.18}
\usetikzlibrary{arrows.meta,calc,positioning,patterns,decorations.pathreplacing,intersections,shapes.misc,backgrounds}

\definecolor{tceemorado}{HTML}{5F2987}
\definecolor{tceedorado}{HTML}{B8860B}
\definecolor{tceeverde}{HTML}{3C7D3A}
\definecolor{tceeazul}{HTML}{1F5F99}
\definecolor{tceerojo}{HTML}{B22222}
\definecolor{tceenaranja}{HTML}{E07B28}
\definecolor{tceegris}{HTML}{767171}
\definecolor{tceesalvia}{HTML}{E2EFD9}

% Capas para el vídeo: \capa{n}{…} dibuja su contenido solo a partir del paso n.
% En el PDF, la web y Word se ve todo; generar-video.py compila el gráfico una
% vez por paso con \def\tceepaso{k} para que cada elemento aparezca al narrarlo.
% Vectores con flecha encima, como en el texto (tcee.sty): \vect{x}, \vect{p}
\providecommand{\vect}[1]{\vec{#1}}
\providecommand{\tceepaso}{1000}
\newcommand{\capa}[2]{\ifnum#1>\tceepaso\relax\else#2\fi}
% Gráficos por pasos del mismo tamaño (propuesta del 3.A.29, 11/11/2026): con
% \def\tceefantasma{1} antes del \input (p. ej., x-paso1.tex:
% «\def\tceepaso{3}\def\tceefantasma{1}\input{x.tex}»), las capas posteriores
% al paso se dibujan invisibles y ocupan su sitio. Sin \tceefantasma (el vídeo),
% no cambia nada.
\ifdefined\tceefantasma
  \renewcommand{\capa}[2]{\ifnum#1>\tceepaso\relax\begin{scope}[opacity=0]#2\end{scope}\else#2\fi}
\fi

\tikzset{
  tcee/.style={line cap=round,line join=round,font=\small,>={Stealth[length=2.2mm]}},
  eje/.style={->,thick,black},
  curva1/.style={very thick,tceemorado},
  curva2/.style={very thick,tceeazul},
  curva3/.style={very thick,tceeverde},
  curva4/.style={very thick,tceedorado},
  curvanueva/.style={very thick,tceerojo},
  desplazada/.style={very thick,dashed},
  guia/.style={thin,dashed,tceegris},
  punto/.style={circle,fill=black,inner sep=1.3pt},
  etiqueta/.style={font=\small,inner sep=2pt},
  flecha/.style={->,thick,tceerojo},
  area/.style={fill=tceemorado,fill opacity=0.18,draw=none},
  area2/.style={fill=tceedorado,fill opacity=0.22,draw=none},
}

% \ejes{etiqueta x}{etiqueta y}{largo x}{alto y}
\newcommand{\ejes}[4]{%
  \draw[eje] (0,0) -- (#3,0) node[below left=2pt and 0pt] {#1};
  \draw[eje] (0,0) -- (0,#4) node[left] {#2};
  \node[below left] at (0,0) {$0$};}
% \proyeccion{(x,y)}{etq x}{etq y}: líneas guía a los ejes con etiquetas
\newcommand{\proyeccion}[3]{%
  \draw[guia] let \p1=#1 in (\x1,0) node[below] {#2} |- (0,\y1) node[left] {#3};}

% --------------------------------------------------------------------------
% Paleta «dictamen» (dic…): gráficos de datos de la Guía del dictamen (1D01) y
% de cualquier gráfico de coyuntura que quiera los mismos colores. Un concepto,
% un color en todos los gráficos (CRITERIOS §7, «tres reglas», 24/10/2026);
% tabla completa en fuentes/1D01/_trabajo/paleta-1D01.md. No cambia nada de
% lo anterior: son colores y estilos nuevos.
% Actividad: el PIB (y sus variantes: PIB por persona, PIB mundial), morado.
\definecolor{dicpib}{HTML}{5F2987}
% Demanda interna, gama cálida: consumo (naranjas: privado oscuro, público
% claro), inversión (verde) y existencias (gris con trama, aparte).
\definecolor{dicconspriv}{HTML}{C76A1F}
\definecolor{dicconspub}{HTML}{F5C690}
\definecolor{dicfbcf}{HTML}{2A6B3C}
\definecolor{dicexistencias}{HTML}{B4B4B4}
% Demanda externa y comercio, gama fría (azul claro, que se separa del morado)
\definecolor{dicexterna}{HTML}{6AA6DA}
% Zonas: zona del euro (azul marino), España (rojo), resto de países (grises)
\definecolor{dicze}{HTML}{17406D}
\definecolor{dicespana}{HTML}{C8382E}
\definecolor{dicotro}{HTML}{8A8A8A}
\definecolor{dicotroclaro}{HTML}{B5B5B5}
\definecolor{dicotromasclaro}{HTML}{DADADA}
% Precios: inflación general (negro, continuo) y subyacente (negro,
% discontinuo); componentes: energía (bermellón), alimentos (amarillo) y los
% de la subyacente en azules (bienes industriales claro, servicios oscuro).
\definecolor{dicinflacion}{HTML}{1A1A1A}
\definecolor{dicenergia}{HTML}{D55E00}
\definecolor{dicalimentos}{HTML}{F2D45C}
\definecolor{dicbine}{HTML}{7FB2E0}
\definecolor{dicservicios}{HTML}{1F4E89}
% Trabajo: empleo (verde azulado), horas (el mismo, claro), tasa de actividad
% (verde azulado oscuro), paro (teja); productividad (dorado de la web).
\definecolor{dicempleo}{HTML}{1B7F6B}
\definecolor{dichoras}{HTML}{86CBB8}
\definecolor{dicactividad}{HTML}{0F5247}
\definecolor{dicparo}{HTML}{A23B2A}
\definecolor{dicproductividad}{HTML}{B8860B}
% Hogares: renta disponible (pizarra) y tasa de ahorro (negro); el consumo es
% dicconspriv.
\definecolor{dicrenta}{HTML}{4A5A6A}
\definecolor{dicahorro}{HTML}{1A1A1A}
% Política monetaria: tipos del BCE en azules (depósito, el de la zona del
% euro; operaciones principales, medio; facilidad marginal, claro); balance:
% préstamos a los bancos (dorado claro), resto del activo (gris claro).
\definecolor{dicdeposito}{HTML}{17406D}
\definecolor{dicopf}{HTML}{4F86C6}
\definecolor{dicfmc}{HTML}{9EC5EA}
\definecolor{dicprestamos}{HTML}{D9A520}
\definecolor{dicresto}{HTML}{CFCFCF}
% Ciclo: series adelantadas, coincidentes (el PIB) y retrasadas
\definecolor{dicadelantada}{HTML}{D55E00}
\definecolor{dicretrasada}{HTML}{6AA6DA}
% Referencias: rejilla, previsión, sombreados
\definecolor{dicreferencia}{HTML}{767171}

\tikzset{
  % barra de un componente: \addplot[dicbarra=dicfbcf]
  dicbarra/.style={fill=#1, draw=#1!65!black, line width=0.3pt},
  % existencias: gris con trama, para que se distinga también en gris
  dicbarratrama/.style={fill=#1, draw=#1!65!black, line width=0.3pt,
    postaction={pattern=north east lines, pattern color=white}},
  % línea de una serie; con halo blanco cuando pasa sobre barras
  diclinea/.style={very thick, #1},
  diclineahalo/.style={very thick, #1, preaction={draw=white, line width=3.4pt, solid}},
  % puntos de una serie (capa aparte, después de la línea)
  dicpuntos/.style={only marks, mark=*, mark size=1.8pt,
    mark options={fill=white, draw=#1, line width=0.9pt, solid}},
  % línea del cero o de un umbral, después de las barras
  diccero/.style={thin, black},
  dicumbral/.style={thin, dashed, dicreferencia},
}
% Capas de los gráficos de datos (CRITERIOS §7, «tres reglas»), de abajo
% arriba: cuadrícula → sombreados (dicfondo) → barras y áreas (dicbarras) → ejes, cero y umbrales
% (axis lines) → líneas (main) → puntos (dicpuntos) → rótulos (axis
% descriptions). No basta el orden del código: pgfplots dibuja cuadrícula,
% ejes y marcas en sus propias capas (y area style pone los ejes encima), y
% dentro de un axis los \draw, \fill y \node no admiten capa (van con los
% datos de main); por eso todo lo que se dibuja dentro del axis es un \addplot.
% Uso: \begin{axis}[…, dicapas] (al final de las opciones);
%   barras y áreas: \addplot[on layer=dicbarras, …];
%   sombreados:     \dicsombra[estilo]{x1}{y1}{x2}{y2} (rectángulo, en dicfondo,
%                   bajo las barras y sobre la cuadrícula);
%   cero y umbrales: \dicraya[estilo]{(x1,y1) (x2,y2)} (en axis lines);
%   líneas: \addplot normal (main); puntos: \addplot[on layer=dicpuntos, …]
%   (las marcas de cualquier serie van solas a dicpuntos);
%   rótulos: \dicrotulo[opciones del nodo]{(x,y)}{texto} (en axis descriptions),
%   siempre sin fondo (nada de fill=white: si chocan con una línea, se desplazan;
%   CRITERIOS §7, regla 4).
% Los \capa{n} del vídeo no cambian: siguen envolviendo instrucciones completas.
\pgfdeclarelayer{dicfondo}
\pgfdeclarelayer{dicbarras}
\pgfdeclarelayer{dicpuntos}
\pgfplotsset{
  % (el segundo argumento reparte cuadrícula, ejes y rótulos como el estándar)
  layers/dic/.define layer set={axis background, axis grid, dicfondo, dicbarras,
    axis ticks, axis lines, axis tick labels, pre main, main, dicpuntos,
    axis descriptions, axis foreground}{/pgfplots/layers/standard},
  % leyenda sin caja blanca: toma el fondo del gráfico, también si el gráfico va
  % dentro de un cuadro gris (CRITERIOS §7, regla 4, 26/10/2026)
  dicapas/.style={axis on top=false, set layers=dic, mark layer=dicpuntos,
    legend style={fill=none}},
  dicsuelto/.style={forget plot, stack plots=false, sharp plot, mark=none},
}
\newcommand{\dicraya}[2][]{%
  \addplot[on layer=axis lines, dicsuelto, diccero, #1] coordinates {#2};}
\newcommand{\dicsombra}[5][]{%
  \addplot[on layer=dicfondo, dicsuelto, draw=none, #1]
    coordinates {(#2,#3) (#4,#3) (#4,#5) (#2,#5)};}
\newcommand{\dicrotulo}[3][]{%
  \addplot[on layer=axis descriptions, dicsuelto, draw=none] coordinates {#2}
    node[pos=0, #1] {#3};}

%   \begin{document}
%   \begin{tikzpicture}[tcee]
%     \ejes{Q}{P}{6}{5}
%     \draw[curva1] (0.5,4.5) -- (5,0.8) node[etiqueta,right] {$D$};
%     \capa{2}{\draw[curva2] ...;}   % en el vídeo aparece en el paso 2
%   \end{tikzpicture}
%   \end{document}
%
% Paleta: la de la web (morado, dorado, salvia) más tres colores de apoyo
% distinguibles en blanco y negro por el trazo.
% ==========================================================================
\usepackage[T1]{fontenc}
\usepackage{newpxtext,newpxmath}
\usepackage{xcolor}
\usepackage{tikz,pgfplots}
\pgfplotsset{compat=1.18}
\usetikzlibrary{arrows.meta,calc,positioning,patterns,decorations.pathreplacing,intersections,shapes.misc,backgrounds}

\definecolor{tceemorado}{HTML}{5F2987}
\definecolor{tceedorado}{HTML}{B8860B}
\definecolor{tceeverde}{HTML}{3C7D3A}
\definecolor{tceeazul}{HTML}{1F5F99}
\definecolor{tceerojo}{HTML}{B22222}
\definecolor{tceenaranja}{HTML}{E07B28}
\definecolor{tceegris}{HTML}{767171}
\definecolor{tceesalvia}{HTML}{E2EFD9}

% Capas para el vídeo: \capa{n}{…} dibuja su contenido solo a partir del paso n.
% En el PDF, la web y Word se ve todo; generar-video.py compila el gráfico una
% vez por paso con \def\tceepaso{k} para que cada elemento aparezca al narrarlo.
% Vectores con flecha encima, como en el texto (tcee.sty): \vect{x}, \vect{p}
\providecommand{\vect}[1]{\vec{#1}}
\providecommand{\tceepaso}{1000}
\newcommand{\capa}[2]{\ifnum#1>\tceepaso\relax\else#2\fi}
% Gráficos por pasos del mismo tamaño (propuesta del 3.A.29, 11/11/2026): con
% \def\tceefantasma{1} antes del \input (p. ej., x-paso1.tex:
% «\def\tceepaso{3}\def\tceefantasma{1}\input{x.tex}»), las capas posteriores
% al paso se dibujan invisibles y ocupan su sitio. Sin \tceefantasma (el vídeo),
% no cambia nada.
\ifdefined\tceefantasma
  \renewcommand{\capa}[2]{\ifnum#1>\tceepaso\relax\begin{scope}[opacity=0]#2\end{scope}\else#2\fi}
\fi

\tikzset{
  tcee/.style={line cap=round,line join=round,font=\small,>={Stealth[length=2.2mm]}},
  eje/.style={->,thick,black},
  curva1/.style={very thick,tceemorado},
  curva2/.style={very thick,tceeazul},
  curva3/.style={very thick,tceeverde},
  curva4/.style={very thick,tceedorado},
  curvanueva/.style={very thick,tceerojo},
  desplazada/.style={very thick,dashed},
  guia/.style={thin,dashed,tceegris},
  punto/.style={circle,fill=black,inner sep=1.3pt},
  etiqueta/.style={font=\small,inner sep=2pt},
  flecha/.style={->,thick,tceerojo},
  area/.style={fill=tceemorado,fill opacity=0.18,draw=none},
  area2/.style={fill=tceedorado,fill opacity=0.22,draw=none},
}

% \ejes{etiqueta x}{etiqueta y}{largo x}{alto y}
\newcommand{\ejes}[4]{%
  \draw[eje] (0,0) -- (#3,0) node[below left=2pt and 0pt] {#1};
  \draw[eje] (0,0) -- (0,#4) node[left] {#2};
  \node[below left] at (0,0) {$0$};}
% \proyeccion{(x,y)}{etq x}{etq y}: líneas guía a los ejes con etiquetas
\newcommand{\proyeccion}[3]{%
  \draw[guia] let \p1=#1 in (\x1,0) node[below] {#2} |- (0,\y1) node[left] {#3};}

% --------------------------------------------------------------------------
% Paleta «dictamen» (dic…): gráficos de datos de la Guía del dictamen (1D01) y
% de cualquier gráfico de coyuntura que quiera los mismos colores. Un concepto,
% un color en todos los gráficos (CRITERIOS §7, «tres reglas», 24/10/2026);
% tabla completa en fuentes/1D01/_trabajo/paleta-1D01.md. No cambia nada de
% lo anterior: son colores y estilos nuevos.
% Actividad: el PIB (y sus variantes: PIB por persona, PIB mundial), morado.
\definecolor{dicpib}{HTML}{5F2987}
% Demanda interna, gama cálida: consumo (naranjas: privado oscuro, público
% claro), inversión (verde) y existencias (gris con trama, aparte).
\definecolor{dicconspriv}{HTML}{C76A1F}
\definecolor{dicconspub}{HTML}{F5C690}
\definecolor{dicfbcf}{HTML}{2A6B3C}
\definecolor{dicexistencias}{HTML}{B4B4B4}
% Demanda externa y comercio, gama fría (azul claro, que se separa del morado)
\definecolor{dicexterna}{HTML}{6AA6DA}
% Zonas: zona del euro (azul marino), España (rojo), resto de países (grises)
\definecolor{dicze}{HTML}{17406D}
\definecolor{dicespana}{HTML}{C8382E}
\definecolor{dicotro}{HTML}{8A8A8A}
\definecolor{dicotroclaro}{HTML}{B5B5B5}
\definecolor{dicotromasclaro}{HTML}{DADADA}
% Precios: inflación general (negro, continuo) y subyacente (negro,
% discontinuo); componentes: energía (bermellón), alimentos (amarillo) y los
% de la subyacente en azules (bienes industriales claro, servicios oscuro).
\definecolor{dicinflacion}{HTML}{1A1A1A}
\definecolor{dicenergia}{HTML}{D55E00}
\definecolor{dicalimentos}{HTML}{F2D45C}
\definecolor{dicbine}{HTML}{7FB2E0}
\definecolor{dicservicios}{HTML}{1F4E89}
% Trabajo: empleo (verde azulado), horas (el mismo, claro), tasa de actividad
% (verde azulado oscuro), paro (teja); productividad (dorado de la web).
\definecolor{dicempleo}{HTML}{1B7F6B}
\definecolor{dichoras}{HTML}{86CBB8}
\definecolor{dicactividad}{HTML}{0F5247}
\definecolor{dicparo}{HTML}{A23B2A}
\definecolor{dicproductividad}{HTML}{B8860B}
% Hogares: renta disponible (pizarra) y tasa de ahorro (negro); el consumo es
% dicconspriv.
\definecolor{dicrenta}{HTML}{4A5A6A}
\definecolor{dicahorro}{HTML}{1A1A1A}
% Política monetaria: tipos del BCE en azules (depósito, el de la zona del
% euro; operaciones principales, medio; facilidad marginal, claro); balance:
% préstamos a los bancos (dorado claro), resto del activo (gris claro).
\definecolor{dicdeposito}{HTML}{17406D}
\definecolor{dicopf}{HTML}{4F86C6}
\definecolor{dicfmc}{HTML}{9EC5EA}
\definecolor{dicprestamos}{HTML}{D9A520}
\definecolor{dicresto}{HTML}{CFCFCF}
% Ciclo: series adelantadas, coincidentes (el PIB) y retrasadas
\definecolor{dicadelantada}{HTML}{D55E00}
\definecolor{dicretrasada}{HTML}{6AA6DA}
% Referencias: rejilla, previsión, sombreados
\definecolor{dicreferencia}{HTML}{767171}

\tikzset{
  % barra de un componente: \addplot[dicbarra=dicfbcf]
  dicbarra/.style={fill=#1, draw=#1!65!black, line width=0.3pt},
  % existencias: gris con trama, para que se distinga también en gris
  dicbarratrama/.style={fill=#1, draw=#1!65!black, line width=0.3pt,
    postaction={pattern=north east lines, pattern color=white}},
  % línea de una serie; con halo blanco cuando pasa sobre barras
  diclinea/.style={very thick, #1},
  diclineahalo/.style={very thick, #1, preaction={draw=white, line width=3.4pt, solid}},
  % puntos de una serie (capa aparte, después de la línea)
  dicpuntos/.style={only marks, mark=*, mark size=1.8pt,
    mark options={fill=white, draw=#1, line width=0.9pt, solid}},
  % línea del cero o de un umbral, después de las barras
  diccero/.style={thin, black},
  dicumbral/.style={thin, dashed, dicreferencia},
}
% Capas de los gráficos de datos (CRITERIOS §7, «tres reglas»), de abajo
% arriba: cuadrícula → sombreados (dicfondo) → barras y áreas (dicbarras) → ejes, cero y umbrales
% (axis lines) → líneas (main) → puntos (dicpuntos) → rótulos (axis
% descriptions). No basta el orden del código: pgfplots dibuja cuadrícula,
% ejes y marcas en sus propias capas (y area style pone los ejes encima), y
% dentro de un axis los \draw, \fill y \node no admiten capa (van con los
% datos de main); por eso todo lo que se dibuja dentro del axis es un \addplot.
% Uso: \begin{axis}[…, dicapas] (al final de las opciones);
%   barras y áreas: \addplot[on layer=dicbarras, …];
%   sombreados:     \dicsombra[estilo]{x1}{y1}{x2}{y2} (rectángulo, en dicfondo,
%                   bajo las barras y sobre la cuadrícula);
%   cero y umbrales: \dicraya[estilo]{(x1,y1) (x2,y2)} (en axis lines);
%   líneas: \addplot normal (main); puntos: \addplot[on layer=dicpuntos, …]
%   (las marcas de cualquier serie van solas a dicpuntos);
%   rótulos: \dicrotulo[opciones del nodo]{(x,y)}{texto} (en axis descriptions),
%   siempre sin fondo (nada de fill=white: si chocan con una línea, se desplazan;
%   CRITERIOS §7, regla 4).
% Los \capa{n} del vídeo no cambian: siguen envolviendo instrucciones completas.
\pgfdeclarelayer{dicfondo}
\pgfdeclarelayer{dicbarras}
\pgfdeclarelayer{dicpuntos}
\pgfplotsset{
  % (el segundo argumento reparte cuadrícula, ejes y rótulos como el estándar)
  layers/dic/.define layer set={axis background, axis grid, dicfondo, dicbarras,
    axis ticks, axis lines, axis tick labels, pre main, main, dicpuntos,
    axis descriptions, axis foreground}{/pgfplots/layers/standard},
  % leyenda sin caja blanca: toma el fondo del gráfico, también si el gráfico va
  % dentro de un cuadro gris (CRITERIOS §7, regla 4, 26/10/2026)
  dicapas/.style={axis on top=false, set layers=dic, mark layer=dicpuntos,
    legend style={fill=none}},
  dicsuelto/.style={forget plot, stack plots=false, sharp plot, mark=none},
}
\newcommand{\dicraya}[2][]{%
  \addplot[on layer=axis lines, dicsuelto, diccero, #1] coordinates {#2};}
\newcommand{\dicsombra}[5][]{%
  \addplot[on layer=dicfondo, dicsuelto, draw=none, #1]
    coordinates {(#2,#3) (#4,#3) (#4,#5) (#2,#5)};}
\newcommand{\dicrotulo}[3][]{%
  \addplot[on layer=axis descriptions, dicsuelto, draw=none] coordinates {#2}
    node[pos=0, #1] {#3};}

\providecommand{\autor}[1]{\textsc{#1}}
\providecommand{\vect}[1]{\vec{#1}}
\begin{document}
\begin{tikzpicture}[tcee,x=1cm,y=1cm,
    banda/.style={fill=#1, text=white, draw=none, rounded corners=3pt, anchor=north west,
      inner sep=4pt, align=left, font=\footnotesize},
    texto/.style={anchor=north west, inner sep=0pt, align=flush left, font=\scriptsize,
      execute at begin node={\setlength{\baselineskip}{8.6pt}}},
    rotfila/.style={anchor=north west, inner sep=0pt, align=flush left,
      font=\scriptsize\bfseries, text=tceemorado}]
  % Columnas: rótulo de la fila; a), b), c) y d)
  \def\xr{0}\def\wr{1.75}
  \def\xa{1.85}\def\wa{3.25}
  \def\xb{5.2}\def\wb{3.45}
  \def\xc{8.75}\def\wc{3.45}
  \def\xd{12.3}\def\wdd{4.35}
  % Filas: arriba (y) y altura (h)
  \def\yH{-0.95}\def\hH{1.4}
  \def\yS{-2.45}\def\hS{3.6}
  \def\yE{-6.15}\def\hE{4.05}
  \def\yP{-10.3}\def\hP{2.05}
  \def\yM{-12.45}\def\hM{2.65}
  \def\yF{-15.2}\def\hF{2.05}
  \def\yL{-17.35}\def\hL{2.35}
  % \celda{x}{w}{y}{h}{color}{texto}
  \newcommand{\celda}[6]{%
    \fill[#5!9, rounded corners=2pt] (#1,#3) rectangle ++(#2,-#4);
    \draw[#5!70, line width=0.5pt, rounded corners=2pt] (#1,#3) rectangle ++(#2,-#4);
    \pgfmathsetmacro{\anchocelda}{#2-0.2}%
    \node[texto, text width=\anchocelda cm] at (#1+0.1,#3-0.1) {#6};}
  \newcommand{\cabecera}[4]{% x, w, color, texto
    \pgfmathsetmacro{\anchocab}{#2-0.28}%
    \node[banda=#3, font=\scriptsize, execute at begin node={\hyphenpenalty=10000\exhyphenpenalty=10000}, text width=\anchocab cm, minimum height=\hH cm, anchor=north west, align=center] at (#1,\yH) {#4};}
  \newcommand{\rotulo}[3]{% y, h, texto
    \node[rotfila, text width=\wr cm] at (\xr,#1-0.1) {#3};}

  % 1: título
  \capa{1}{\node[banda=tceemorado, text width=16.4cm] at (0,0)
    {\textbf{Los 4 modelos del tema, frente a frente}: de los estáticos, con efectos nítidos dado el régimen, a los dinámicos, con efectos que dependen de la economía y de la regla.};}

  % 2: rótulos de las filas
  \capa{2}{\rotulo{\yS}{\hS}{Supuestos clave};}
  \capa{2}{\rotulo{\yE}{\hE}{Ecuaciones};}
  \capa{2}{\rotulo{\yP}{\hP}{Qué determina los precios y el tipo de cambio};}
  \capa{2}{\rotulo{\yM}{\hM}{Política monetaria};}
  \capa{2}{\rotulo{\yF}{\hF}{Política fiscal};}
  \capa{2}{\rotulo{\yL}{\hL}{Límites};}

  % 3: a) IS-LM con OA-DA (economía cerrada)
  \capa{3}{\cabecera{\xa}{\wa}{tceedorado}{\textbf{a) IS-LM con OA-DA}\\ economía cerrada};}
  \capa{3}{\celda{\xa}{\wa}{\yS}{\hS}{tceedorado}{%
    Economía cerrada.\\
    IS-LM: precios fijos.\\
    DA-OA: $P$ se ajusta; salarios nominales que se revisan con retraso.\\
    Expectativas adaptativas o estáticas; sin microfundamentos.}}
  \capa{3}{\celda{\xa}{\wa}{\yE}{\hE}{tceedorado}{%
    IS: $Y=\dfrac{\overline{A}-b\,i}{1-c\,(1-t)}$\\[2pt]
    LM: $\dfrac{M}{P}=k\,Y-h\,i$\\[2pt]
    DA: un equilibrio IS-LM por cada $P$ (efecto \autor{Keynes}).\\[2pt]
    OA: $\pi_t-\pi^e_t=-a\,(U_t-U^n)+\varepsilon_t$; a medio plazo, vertical en $Y^n$.}}
  \capa{3}{\celda{\xa}{\wa}{\yP}{\hP}{tceedorado}{%
    $P$: el corte de la DA y la OA; a medio plazo, la DA sobre la OA vertical.\\
    No hay tipo de cambio.}}
  \capa{3}{\celda{\xa}{\wa}{\yM}{\hM}{tceedorado}{%
    Expansiva: baja $i$, sube la inversión; suben $Y$ y $P$.\\
    A medio plazo, neutral: solo sube $P$.}}
  \capa{3}{\celda{\xa}{\wa}{\yF}{\hF}{tceedorado}{%
    Expansiva: suben $Y$ y $P$; sube $i$ y expulsa inversión.\\
    A medio plazo, solo $P$ y la composición del gasto.}}
  \capa{3}{\celda{\xa}{\wa}{\yL}{\hL}{tceedorado}{%
    Sin comercio ni tipo de cambio; sin expectativas racionales (crítica de \autor{Lucas}) ni criterio de bienestar.}}

  % 4: b) IS-LM-BP con OA-DA (Mundell-Fleming)
  \capa{4}{\cabecera{\xb}{\wb}{tceeazul}{\textbf{b) IS-LM-BP con OA-DA}\\ \autor{Mundell}-\autor{Fleming}, abierta};}
  \capa{4}{\celda{\xb}{\wb}{\yS}{\hS}{tceeazul}{%
    País pequeño: $Y^*$ e $i^*$ dados.\\
    Movilidad del capital $\beta_K$; \autor{Marshall}-\autor{Lerner}; expectativas estáticas.\\
    Precios fijos en \autor{Mundell}-\autor{Fleming}; en la DA-OA, $P$ se ajusta y $P^*$ está dado.\\
    Salarios con el índice de consumo (apertura $\alpha$).}}
  \capa{4}{\celda{\xb}{\wb}{\yE}{\hE}{tceeazul}{%
    IS: $Y=\dfrac{\overline{A}-b\,i+x\,Y^*+\nu\,q}{1-c\,(1-t)+m}$\\[2pt]
    LM: $\dfrac{M}{P}=k\,Y-h\,i$\\[2pt]
    BP: $x\,Y^*-m\,Y+\nu\,q+\beta_K\,(i-i^*)=0$\\[2pt]
    DA: fijo, efecto competitividad; flexible, efecto saldos reales.\\[2pt]
    OA con la cuña: $F(U^n,Z)=\dfrac{1}{(1+\mu)\,q^{\alpha}}$}}
  \capa{4}{\celda{\xb}{\wb}{\yP}{\hP}{tceeazul}{%
    Tipo fijo: se ajusta $M$ (endógeno). Flexible: se ajusta $E$.\\
    $P$: DA-OA. Inflación de consumo: $\pi\simeq(1-\alpha)\,\pi_H+\alpha\,(\gamma_E+\gamma_{P^*})$.}}
  \capa{4}{\celda{\xb}{\wb}{\yM}{\hM}{tceeazul}{%
    Movilidad perfecta: con tipo fijo, \textcolor{tceerojo}{ineficaz} (se van las reservas); con flexible, \textcolor{tceeverde}{muy eficaz} (depreciación).\\
    A medio plazo, neutral.}}
  \capa{4}{\celda{\xb}{\wb}{\yF}{\hF}{tceeazul}{%
    Movilidad perfecta: con tipo fijo, \textcolor{tceeverde}{muy eficaz}; con flexible, \textcolor{tceerojo}{ineficaz} (expulsa $XN$).\\
    A medio plazo, con tipo fijo, la expulsión llega por los precios.}}
  \capa{4}{\celda{\xb}{\wb}{\yL}{\hL}{tceeazul}{%
    Estático; flujos y no fondos; sin restricción intertemporal ni criterio de bienestar.}}

  % 5: c) el modelo neokeynesiano canónico
  \capa{5}{\cabecera{\xc}{\wc}{tceemorado}{\textbf{c) Neokeynesiano canónico}\\ dinámico, cerrado};}
  \capa{5}{\celda{\xc}{\wc}{\yS}{\hS}{tceemorado}{%
    Hogares que optimizan, con expectativas racionales ($\theta$, $\varphi$).\\
    Competencia monopolística ($\sigma$) y precios à la \autor{Calvo} ($\eta$).\\
    Sin dinero (\emph{cashless}); un banco central con una regla de tipos.\\
    Ecuaciones log-linealizadas.}}
  \capa{5}{\celda{\xc}{\wc}{\yE}{\hE}{tceemorado}{%
    IS dinámica:\\ $x_t=\mathbb{E}_t[x_{t+1}]-\frac{1}{\theta}\big(i_t-\mathbb{E}_t[\pi_{t+1}]-r^n_t\big)$\\[2pt]
    NKPC:\\ $\pi_t=\beta\,\mathbb{E}_t[\pi_{t+1}]+\kappa\,x_t+u_t$\\[2pt]
    Regla de \autor{Taylor}:\\ $i_t=\rho+\phi_\pi\,\pi_t+\phi_x\,x_t+\varepsilon_t$\\[2pt]
    Forma reducida:\\ $x_t=\dfrac{\hat r^n_t-\phi_\pi\,u_t-\varepsilon_t}{\theta+\phi_x+\phi_\pi\,\kappa}$}}
  \capa{5}{\celda{\xc}{\wc}{\yP}{\hP}{tceemorado}{%
    La inflación, no el nivel: la NKPC, con las expectativas y la brecha; la regla la ancla.\\
    No hay tipo de cambio.}}
  \capa{5}{\celda{\xc}{\wc}{\yM}{\hM}{tceemorado}{%
    Contractiva ($\varepsilon_t>0$): caen la brecha y la inflación.\\
    Demanda: divina coincidencia; costes: dilema.}}
  \capa{5}{\celda{\xc}{\wc}{\yF}{\hF}{tceemorado}{%
    Entra por $r^n_t$. El multiplicador es menor cuanto más agresiva es la regla; puede superar 1 en el límite inferior de los tipos.}}
  \capa{5}{\celda{\xc}{\wc}{\yL}{\hL}{tceemorado}{%
    Necesita inercia añadida; $u_t$, casi un residuo; sin sector financiero ni tipo de cambio.}}

  % 6: d) la nueva macroeconomía abierta
  \capa{6}{\cabecera{\xd}{\wdd}{tceeverde}{\textbf{d) Nueva macroeconomía abierta}\\ PEA: \autor{Galí}-\autor{Monacelli} y \autor{Furlanetto}\\ 2PG: \autor{Corsetti}-\autor{Pesenti}};}
  \capa{6}{\celda{\xd}{\wdd}{\yS}{\hS}{tceeverde}{%
    \textbf{PEA}: continuo de economías pequeñas; apertura $\alpha$ y sesgo nacional; precios en la moneda del productor; mercados completos; gasto en bienes nacionales, hogares ricardianos.\\[1pt]
    \textbf{2PG}: 2 países, A y B; $U=\ln C-\chi\,\ell$; precios fijados un periodo antes, en la moneda del productor o del comprador; sin gasto público; $\mu$, el tono monetario (no el margen).}}
  \capa{6}{\celda{\xd}{\wdd}{\yE}{\hE}{tceeverde}{%
    \textbf{PEA}: $\pi_t=\pi_{H,t}+\alpha\,\Delta\tau_t$\\
    $x_t=\mathbb{E}_t[x_{t+1}]-\frac{1}{\theta_{\alpha G}}\big(i_t-\mathbb{E}_t[\pi_{H,t+1}]-r^n_t\big)$\\
    $\pi_{H,t}=\beta\,\mathbb{E}_t[\pi_{H,t+1}]+\kappa_\alpha\,x_t+u_t$\\
    $i_t=\rho+\phi_\pi\,\pi_{H,t}+\phi_x\,x_t+\varepsilon_t$, con $\theta_{\alpha G}=\dfrac{\theta}{(1+\alpha\,\Theta)(1-s_G)}$\\[3pt]
    \textbf{2PG}: AD: $C=\dfrac{\mu}{P}$;\quad AS: $C=Z\,\ell\,\tau$;\\[1pt]
    NR: $\overline{\ell}=\dfrac{\sigma-1}{\sigma\,\chi}$;\quad $E=\dfrac{\mu}{\mu^*}$}}
  \capa{6}{\celda{\xd}{\wdd}{\yP}{\hP}{tceeverde}{%
    \textbf{PEA}: $\pi_H$ por la NKPC abierta; $E$ por la paridad de intereses y el reparto del riesgo, $q_t=(1-\alpha)\,\tau_t$; $\pi$ suma lo importado.\\
    \textbf{2PG}: precios fijados por adelantado; $E$, el cociente de los tonos monetarios.}}
  \capa{6}{\celda{\xd}{\wdd}{\yM}{\hM}{tceeverde}{%
    \textbf{PEA}: contractiva: caen la brecha y $\pi_H$, la moneda se aprecia y $\pi$ cae más; la apertura amplifica si $\Theta>0$. Mejor regla: inflación interna $>$ consumo $>$ fijo.\\
    \textbf{2PG}: una expansión en A ayuda a B con precios del productor; con los del comprador, lo empobrece.}}
  \capa{6}{\celda{\xd}{\wdd}{\yF}{\hF}{tceeverde}{%
    \textbf{PEA}: $\dfrac{\partial Y}{\partial G}=\dfrac{\theta_{\alpha G}}{\theta_{\alpha G}+\varphi}<1$, menor cuanto más abierta; aprecia y caen $C$ y $XN$; con tipo fijo, unos 0,9 frente a 0,35.\\
    \textbf{2PG}: apenas desarrollada.}}
  \capa{6}{\celda{\xd}{\wdd}{\yL}{\hL}{tceeverde}{%
    \textbf{PEA}: mercados completos y traspaso completo; el rompecabezas del consumo; los efectos dependen de los parámetros y de la regla.\\
    \textbf{2PG}: depende de la moneda de facturación y de las preferencias logarítmicas.}}

  % 7: la lectura de la tabla
  \capa{7}{\node[banda=tceemorado, text width=16.4cm] at (0,-19.85)
    {\textbf{Cómo se lee}: de izquierda a derecha, cada modelo levanta un supuesto del anterior (los precios fijos, la economía cerrada, las expectativas estáticas, el país aislado). Los de a) y b) dan efectos de signo conocido; los de c) y d) recuperan sus direcciones básicas, pero el tamaño, y a veces el signo, depende de la apertura, las elasticidades, las rigideces, la regla monetaria y la moneda de facturación.};}
\end{tikzpicture}
\end{document}
